\documentclass[a4paper,11pt]{article}

%%%%%%%% CREATE DOCUMENT STRUCTURE %%%%%%%%
%% Language and font encodings
\usepackage[italian]{babel}
\usepackage[utf8x]{inputenc}
\usepackage[T1]{fontenc}

%% Sets page size and margins
\usepackage[a4paper,top=3cm,bottom=2cm,left=2cm,right=2cm,marginparwidth=1.75cm]{geometry}

%% Useful packages
\usepackage{amsmath}
\usepackage{amssymb}
\usepackage{amsfonts} % for /mathbb and similar math symbols
\usepackage{graphicx}
\usepackage{caption}
\usepackage{subcaption}
\usepackage{float}
\usepackage{titling}
\usepackage{blindtext}
\usepackage[square,sort,comma,numbers]{natbib}
\usepackage{xcolor}
\usepackage[colorlinks=true, allcolors=blue, pdfpagelabels]{hyperref}

\newcommand{\HRule}{\rule{\linewidth}{0.5mm}} 	% horizontal line and its thickness

\begin{document}
	%%%% Title Page
	\hypersetup{pageanchor=false}
	\begin{titlepage}
		\begin{center}
			% Some logo, optional
			\includegraphics[width=0.5\textwidth]{img/unipi-logo.png}\\[1cm]
			
			% University
			\textsc{\Large Università di Pisa, Dipartimento di Informatica}\\[1cm]
			
			% Document info
			\textsc{\Large Laboratorio di Reti - A.A. 2021-2022 - corso B}\\[0.2cm]
			\textsc{\large Docente: Federica Paganelli}\\[1cm]
			
			% Assignment title (enclosed between horizontal lines
			\HRule \\[0.8cm]
			{ \Large \bfseries Relazione progetto Winsome}\\[0.7cm]
			\HRule \\[2cm]
			
			% Author and date
			\Large Nicola Vetrini, matricola 600199\\[0.2cm]
			n.vetrini@studenti.unipi.it\\[1cm]
			{\large \today}\\[5cm]
			
			\vfill
		\end{center}
	\end{titlepage}
	\hypersetup{pageanchor=true}
	
	\clearpage
	\tableofcontents
	\clearpage
	
	\section{Introduzione}
	Lo scopo del progetto è stata la realizzazione di una rete sociale, Winsome, basata sulla condivisione	di contenuti tra un insieme di utenti. \\
	Il social network Winsome consiste di due componenti software: \textbf{WinsomeClient} (client nel seguito)
	e \textbf{WinsomeServer} (server nel seguito). Il client ed il server comunicano principalmente tramite un 
	protocollo richiesta/risposta su un socket TCP subito prima della login di ciascun utente e chiuso al logout. Pertanto il client ed il server possono essere avviati in modo indipendente, avendo però cura di notare che molti comandi del client non
	danno alcun risultato se la connessione al server non può essere stabilita.\\
	Vi sono inoltre gli oggetti remoti per la comunicazione tramite RMI e RMI callback: il server crea un registry RMI che utilizza per fornire ai client gli stub per la registrazione e la sottoscrizione al servizio di callback.\\
	Il server può inoltre fornire ai client che lo richiedano l'indirizzo di un gruppo UDP multicast, sul quale notifica regolarmente che i wallet degli utenti sono stati
	aggiornati (non viene comunicata alcuna informazione riguardo le ricompense calcolate). I client eseguono la richiesta di tale indirizzo ed effettuano la join
	su di esso soltanto al login di un utente, e lasciano il gruppo al logout: ho ritenuto non efficiente mantenere un socket in attesa di datagrammi (in un thread separato, come illustrato nel seguito) se non vi è alcun utente che possa controllare ontestualmente lo stato del proprio wallet dopo aver ricevuto la notifica.\\
	
	Un aspetto dell'implementazione di Winsome proposta che assume particolare rilievo è la gestione dei rewin: l'azione compiuta dal rewin di un post è quella di
	spostare un post dal proprio feed al proprio blog. 
	Dato che nella mia implementazione vi è una mappa globale dei post vi erano due approcci possibili: un primo approccio consiste nell'inserire il riferimento al post
	nel blog, l'altro nel creare un nuovo post che contenga l'id del post originale ed il suo autore, con una flag che indica che quel post è un rewin.\\
	Nell'implementazione è stato adottato il secondo approccio, poiché consentirebbe, anche se non è stato implementato, di attribuire una ricompensa anche agli utenti
	che contribuiscono alla diffusione di un post, e quindi al suo successo nella rete sociale. Un punto importante è l'inserimento di voti e commenti: tali modifiche
	al post devono essere effettuate sul post originale, in modo tale che il calcolo dei commenti e dei voti possa correttamente assegnare all'autore, e non a chi ha effettuato
	il rewin, la somma dovuta. Tale decisione è motivata principalmente dal seguente scenario: se vi fossero due blog, A che ha un numero esiguo di followers e B che ha un gran
	numero di followers, ed il blog B effettuasse il rewin di un post di A esso potrebbe ricevere molti commenti e voti, in quanto compare sul feed di un gran numero di blog.
	Di conseguenza il blog B riceverebbe una grande ricompensa, sfruttanto il post prodotto da un altro utente. La logica implementata previene questo tipo di situazioni, assegnando al creatore del post orignale la ricompensa dovuta alla diffusione del post.\\
	Se il post originale al quale un rewin si riferisce venisse cancellato, allora fallirebbe anche l'aggiunta di voti o commenti (il modo in cui questo controllo viene effettuato
	è commentato nel codice), seppure l'id del post cancellato rimanga visibile nella rete sociale sui blog di utenti che hanno effettuato il rewin di tale post.
	La cancellazione in cascata di tale id sarebbe un'operazione in generale molto dispendiosa, che può essere evitata effettuando dei controlli aggiuntivi molto più veloci.\\
	
	\section{Architettura del sistema}
Il social network Winsome consiste di due componenti software: WinsomeClient (client nel seguito)
 e WinsomeServer (server nel seguito). Il client ed il server comunicano principalmente 
 attraverso delle richieste, contenute nel package WinsomeRequests, di vario tipo. 
 Ogni richiesta è convertita in una stringa JSON tramite la libreria Jackson, e da 
 tale stringa si ottiene un arrray di byte, il quale viene trasmesso su un socket TCP che connette
 il client ed il server.\\
 Attraverso il client è possibile inoltre registrare un nuovo utente su Winsome: tale operazione 
 viene eseguita utilizzando RMI, come da specifica (l'iterfaccia implementata dal server è Signup, 
 con il metodo register).\\
 Il comando \verb|list followers| del client non genera una richiesta sincrona al server, ma piuttosto
 è presente una struttura dati nel client il cui contenuto viene restituito. Tale elenco di followers
 viene aggiornato periodicamente dal server utilizzando il meccanismo delle RMI callback:
 il client al login di un utente si registra presso il server (FollowerUpdaterService, metodo \verb|subscribe()|)
TODO: UDP multicast per wallet
	
	\section{Server Winsome}
Il server è contenuto in un package, \textbf{WinsomeServer}, che contiene le seguenti classi principali:
\begin{itemize}
	\item \verb|ServerMain.java|, contenente la classe main del server
	\item \verb|WinsomeServer.java|, contenente la classe che incapsula i dati della rete 
	sociale e le funzionalità principali
	\item \verb|User.java, Post.java, Comment.java| ed altre, che definiscono il formato dei dati
	all'interno della rete sociale
\end{itemize}

\subsection{Argomenti da riga di comando}
Il server può ricevere uno o più dei seguenti parametri da riga di comando, che ne influenzano il comportamento all'avvio.
Si noti che i valori passati da riga di comando hanno la precedenza su quelli eventualmente letti dal file di configurazione caricato.\\
Se non è stato specificato alcun valore per un parametro, verrà usato il valore di default statico presente nel file \verb|ServerConfig.java|.\\
Per un riepilogo delle opzioni disponibili è sufficiente passare la flag \verb|-h| o \verb|--help|, in ogni caso le opzioni
disponibili sono le seguenti:

\begin{itemize}
	\item \verb|-c --configure <FILE>|: path del file di configurazione che il server deve caricare
	\item \verb|-h --help|: messaggio di uso del server
	\item \verb|-p --socket-port <PORT>|: porta sulla quale, se possibile, viene creata la ServerSocketChannel 
	sulla quale il server si mette in ascolto per accettare le connessioni dai client
	\item \verb|-r --registry <PORT>|: porta sulla quale, se possibile, viene creato il registry 
	RMI utilizzato sia per la registrazione, che per l'iscrizione/disiscrizione dal servizio di aggiornamento dei followers
\end{itemize}

\subsection{Configurazione}
Il server è configurato attraverso un file JSON che viene cercato, nell'ordine, ai seguenti path:
\begin{enumerate}
	\item il path passato attraverso l'opzione \verb|-c <path>|
	\item \verb|config.json| nella directory corrente
	\item \verb|data/WinsomeServer/config.json| il file di configurazione di default
\end{enumerate}
Se tutte le opzioni precedenti non contengono un file di configurazione valido allora il server termina immediatamente.\\
I parametri configurabili sono i seguenti:
\begin{description}
	\item[dataDir]: Il path alla directory contenente i dati del server, ovvero la directory
	che contiene il file \verb|users.json|, la directory \verb|blogs|, etc ...
	\item[registryPort]: La porta del registry RMI nel quale inserire lo stub per la registrazione. Il valore deve essere un intero nel range [0, 65535]
	\item[serverSocketAddress]: Nome host o indirizzo IP del ServerSocket creato dal server per accettare le connessioni dai client (IP pubblico del server nel caso generale, ma in questo caso localhost)
	\item[serverSocketPort]: Porta alla quale deve essere legato il socket sopracitato, nel range [0, 65535]
	\item[multicastGroupAddress]: Indirizzo IP del gruppo multicast sul quale inviare le notifiche di aggiornamento dei wallet
	\item[multicastGroupPort]: Porta relativa all'indirizzo multicast precedente
	\item[minPoolSize]: Intero ($>0$) che rappresenta il numero di core thread nella threadpool per la gestione delle richieste del server
	\item[maxPoolSize]: Intero ($x : INT\_MIN \le minPoolSize \le x < INT\_MAX$) che rappresenta il numero massimo di thread che possono essere gestiti contemporaneamente dalla threadpool
	\item[workQueueSize]: Dimensione della coda di task che la threadpool può accumulare in attesa che un thread del pool sia libero, in accordo con la politica di gestione delineata dalla documentazione di ThreadPoolExecutor
	\item[retryTimeout]: long ($>0$) che rappresenta il numero di millisecondi che l'handler per la gestione delle richieste rifiutate dal threadpool attende prima di provare a sottomettere nuovamente la richesta
	\item[callbackInterval]: long ($>0$) che rappresenta l'intervallo tra due consecutivi aggiornamenti alla lista dei follower del client registrato al servizio. L'unità di misura che quantifica il valore è specificata dal parametro indicato di seguito
	\item[callbackIntervalUnit]: Unità di misura di callbackInterval. Essendo rappresentato con una\\
	\hyperref{https://docs.oracle.com/en/java/javase/11/docs/api/java.base/java/util/concurrent/TimeUnit.html#enum.constant.summary}{category}{name}{TimeUnit} i valori possibili sono quelli dell'enumerazione indicati nel link. Di default l'unità di tempo sono i secondi. Nel file JSON una TimeUnit è serializzata come stringa, quindi per indicare di misurare callbackInterval in secondi è sufficiente scrivere
	"callbackIntervalUnit": "SECONDS" nel file JSON.
	\item[authorPercentage]: Percentuale della ricompensa calcolata che deve essere assegnata all'autore. 
	Il valore specificato deve essere un double $x : 0 < x < 1$
\end{description}
Non è necessario specificare tutti i parametri nel file, poiché quelli assenti assumeranno i valori di default definiti nella classe \verb|ServerConfig.java|. Il file JSON viene deserializzato utilizzando Jackson con ObjectMappper in un'istanza della classe menzionata, ed un riferimento ad essa viene passato alla creazione del server.\\
Nella classe è in realtà presente anche il campo (non serializzato) configFile, che serve soltanto per determinare, nella funzione \verb|getServerConfiguration()| in \verb|main()|, se era stato settato con l'opzione -c il path per un file di configurazione, che ha la precedenza rispetto a quello di default.

\subsection{Esecuzione}
Il codice del server Winsome è contenuto all'interno del package WinsomeServer; la classe ServerMain in tale package contiene il metodo main, quindi è quella da eseguire per far partire il server. Il server stampa su standard output all'avvio tutti i parametri della propria configurazione, con altre informazioni di supporto relative al caricamento
dei dati ed eventuali errori.\\
In particolare, prima viene caricato il file degli utenti; al termine, per ciascun utente, viene caricato il suo blog. Se il file contenente il blog 
non esiste viene notificato su stderr ed il caricamento prosegue (si crea nel server un nuovo blog vuoto per tale utente). Successivamente viene creato il ServerSocketChannel
non bloccante che si mette in ascolto di nuove connessioni e viene attivato un selector per il multiplexing delle richieste sui socket attivi.\\
Quando un client richiede la connessione viene creato il SocketChannel non bloccante e viene registrato per l'operazione di lettura, con l'attachment \textbf{ClientData} che incapsula lo stato di quel client. Quando un SocketChannel diventa readable si effettua la read della richiesta e si registra quel SocketChannel sono per l'operazione di scrittura (\verb|setWritable()|), quindi il selector non è più in
ascolto per operazioni di lettura da quel socket). Una volta pronta la risposta ed il socket per la scrittura la si effettua; successivamente si registra il socket solo per la lettura (\verb|setReadable()|) in modo analogo.\\
La terminazione del server avviene interrompendo la JVM, ad esempio con \verb|Ctrl+C|, ma  dato che all'avvio sono registrati degli shutdown hook per la sincronizzazione dei contenuti del social network lo stato corrente al momento della terminazione viene salvato sovrascrivendo il file degli utenti, dei blog ed il file di configurazione. Una volta terminati i thread di sincronizzazione (tutti operanti su un set di dati indipendenti, quindi non soggetti a particolare attenzione per quanto riguarda la sincronizzazione) la JVM su cui eseguiva il server può terminare.
\subsubsection{Thread e gestione della concorrenza}
Il server Winsome è multithreaded: vi è il thread del main, il cui compito è l'inizializzazione del server e del registry; dal thread principale è fatta partire un istanza di \verb|WinsomeServer|, sottoclasse di \verb|Thread|, che si occupa di gestire lo smistamento delle richieste. Una volta lanciato con successo il WinsomeServer il thread main termina.\\
Nel thread del WinsomeServer vi è un selector che permette di leggere le richieste provenienti dai client e sottometterle ad una threadpool le cui dimensioni minime e massime sono fissate dal file di configurazione.\\
Nel costruttore di WinsomeServer, inoltre, viene attivato un thread per leggere gli utenti di Winsome dal file \verb|users.json| e caricarli in memoria. Una volta letti gli utenti viene attivato un \verb|BlogLoaderThread| per ciascun utente letto dal file, il cui compito è deserializzare dal file \verb|<datadir>/blogs/<user>.json| tutti i post presenti in tale blog e caricarli nelle strutture dati del server (ogni post è inserito nella mappa post dei globale e nella lista dei post di ogni blog). Tali thread terminano una volta caricato il blog (il thread che esegue il costruttore di WinsomeServer si blocca finché la join su ciascuno di essi non ritorna).\\
Alla terminazione del server, poiché lo stato deve persistere, ogni utente ed ogni post dei loro blog devono essere scritti su file. Per fare questo è stato utilizzato
il meccanismo degli shutdown hook, settati come prima istruzione del metodo \verb|run()| del WinsomeServer.\\
Vi è un hook (thread su cui non è stato invocato \verb|start()|) per la sincronizzazione del file degli utenti: \textbf{SyncUsersThread.java}. L'altro hook è \textbf{SyncBlogsThread.java}, il cui unico compito è quello di creare e far partire un \textbf{SyncPostsThread.java} per ogni utente Winsome, che sincronizza i post di un solo blog. Tali thread operano su dati totalmente indipendenti, per cui non è richiesta alcuna sincronizzazione delle loro operazioni, il che consente di avere il massimo grado di parallelizzazione del processo consentito dalla macchina. L'unica accortezza è che il thread che li ha creati non termini fino a che tutti questi thread non sono terminati, altrimenti la JVM potrebbe terminare lasciando uno o più blog in uno stato inconsistente, che provovcherebbe degli errori al riavvio del server.

\subsubsection{Strutture dati}
	
	\section{Client Winsome}
client
	
	\section{Compilazione ed esecuzione}
I comandi elencati di seguito assumono che la directory corrente sia quella base del progetto
(quella con la subdirectory src). Sono inoltre allegati un Makefile e degli script per la compilazione ed esecuzione (\textbf{serv\_run.sh} e \textbf{client\_run.sh}), il cui scopo è semplicemente quello di automatizzare l'esecuzione dei comandi riportati di seguito.\\
La compilazione ed esecuzione è stata testata principalmente con Java 11 su Linux, ma dovrebbe compilare ed eseguire senza modifiche su Java >= 8.\\
Una nota di rilievo: nei comandi elencati di seguito, come in quelli presenti nel Makefile, il carattere separatore di path nella flag \verb|-cp| è ":". Se l'esecuzione viene testata in un sistema Windows è necessario sostituire le occorrenze di ":" con ";".

\subsection{Compilazione}
Il comando per la compilazione del server è il seguente:\\
\verb|javac -d bin -cp "libs/*" -sourcepath src/ src/Winsome/WinsomeServer/*.java|\\
Per compilare il client è sufficiente sostituire WinsomeClient al posto di WinsomeServer.\\
Il risultato della compilazione è nella directory \verb|bin/Winsome|.\\

Sono disponibili anche i comandi per la generazione di eseguibili in formato jar come segue:\\
\verb|jar cfm bin/WinsomeServer.jar server-manifest.txt -C bin Winsome| per il server e\\
\verb|jar cfm bin/WinsomeClient.jar client-manifest.txt -C bin Winsome| per il client.

\subsection{Esecuzione}
Per eseguire il server è sufficiente il comando\\
\verb|java -cp ".:bin/:libs/*" Winsome.WinsomeServer.ServerMain|\\
Sostituendo a "WinsomeServer.ServerMain" "WinsomeClient.ClientMain" si può eseguire il client.\\
Eventuali argomenti da riga di comando al client o al server possono essere aggiunti in fondo.\\
Per eseguire i jar è sufficiente il comando \verb|java -jar bin/WinsomeServer.jar| ed uno analogo per il client.
Si noti che se gli eseguibili sono spostati in una directory diversa non troveranno le librerie esterne utilizzate,
a meno di non ricreare gli archivi in modo opportuno (sostanzialmente modificando \verb|client-manifest.txt| e
 \verb|server-manifest.txt| ed eseguendo \verb|make jars|).

\subsection{Documentazione} \label{javadoc}
È disponibile nel Makefile un target (\verb|make doc|) per generare con javadoc la documentazione delle classi realizzate, navigabile tramite browser.
Il comando (eliminado gli \verb|\|) può essere eseguito anche da terminale. La documentazione viene generata nella subdirectory doc, per cui è sufficiente aprire il file \verb|doc/index.html| con il browser per sfogliarla.\\
Con l'opzione -link di javadoc si ottengono link cliccabili alla documentazione delle classi della libreria standard Java e delle librerie utilizzate, anche se ciò potrebbe introdurre un leggero ritardo nella generazione, per cui le righe contenenti -link possono essere tolte, se necessario.

\subsection{Testing}
Sono incluse nel codice consegnato le directory data e testcases. In data è stata inclusa una rete sociale di esempio, con alcuni post e commenti. Nelle directory testcases vi sono
le reti sociali create per verificare i testcase forniti nella specifica per la formula del calcolo delle ricompense.\\
Un parametro che tipicamente sarà necessario modificare è netIf nel file di configurazione del client, da sostituire con il nome di una interfaccia di rete presente sulla propria
macchina, che consenta l'invio e la ricezione di datagrammi UDP su gruppi multicast. L'altro parametro riguarda la frequenza degli update della lista di followers e del calcolo delle ricompense, settata ad un intervallo breve (20 o 30 secondi) per il testing, che può interferire con l'output del client (avrei potuto scrivere su uno stream diverso da stdout, ma avrei dovuto implementare una qualche forma di interfaccia grafica).
	
	\section{Librerie utilizzate}
Le librerie esterne utilizzate dal software sono state incluse, sotto forma di file .jar, nella
cartella \textit{libs}

\subsection{Apache Commons CLI}
È stata utilizzata la libreria Apache Commons CLI versione 1.5.0 per il parsing degli argomenti da riga 
di comando sia in WinsomeServer che in WinsomeClient.
\subsection{Jackson}
All'interno del software, soprattutto nella componente server, è stata utilizzata ripetutamente
la serializzazione e deserializzazione tra classi ed oggetti JSON scritti su file.
A tale scopo ho scelto di utilizzare la libreria Jackson, versione 2.9.7.
\end{document}